\scanpage{9}
Si on étend les scalaires du plan \inserted{vect} réel
\unsure{au|en} plan vectoriel complexe, l'hom \inserted{$u_{\zeta}$} de
multiplication complexe par $\zeta$ (pour tt nb
complexe) se décompose, suivant deux
droites ``isotropes'' canoniques en
la matrice \inserted{diagonale} $(\zeta,\zeta^{-1})$,
\diagram{\begin{tikzcd}[ampersand replacement=\&,row sep=small,column sep=small]
E\arrow[d,dash]\arrow[dr,shift right=2pt,"\varphi_1"']\arrow[dr,shift left=2pt,"\varphi_2"]\&{}\\
\mathbb{R}\arrow[r]\&\mathbb{C}
\end{tikzcd}}
\[
\left\{\begin{aligned}
\varphi_1(z)\deleted{\unsure{\mathbb C=\mathrm{id}_{\mathbb C}|z=\mathrm{id}_{\mathbb C}}}&=z\\
\varphi_2(z)&=\bar z
\end{aligned}\right.
\]
\[
E_{\mathbb C}\simeq\mathbb C^2
\]
le pt 1 de $E$ devient
le pt $s_0=(1,1)$ de $\mathbb C^2$,
et $s_i$ devient le pt de coordonnées
complexes $(\zeta^i,\zeta^{-i})=u_\zeta^i(s_0)$. La symétrie
p.r. au sommet $s_0$, i.e. $z\longrightarrow\bar z$, est
transformé en l'autom $(\deleted{\unsure{-|x}}x,y)\longmapsto(y,x)$
symétrie des deux axes. On verra que
cette structure du $n$-polygône régulier
est typique, pour les corps \inserted{alg clos} de car.
première à $n$\,; et pour ceux-ci on a la $=$
classification des $n$-polygônes réguliers \inserted{$((n,p)=1)$} que
celle indiquée dans le cas du corps des
réels. (NB Le fait que dans le cas
particulier $k=\mathbb R$, \uline{bien que $k$ ne soit pas
alg} clos, on ait la $=$ classification, est
remarquable et demande de l'attention\ldots
% archive sheet number: 8

